\newpage

\section*{Capítulo 3. Delimitación computacional del problema}
\addcontentsline{toc}{section}{Capítulo 3. Delimitación computacional del problema}
\setcounter{section}{3}
\setcounter{subsection}{0}

En este capítulo se delimita el problema desde una perspectiva computacional, considerando la naturaleza visual de los datos de entrada, la visión por computadora como área de estudio y la justificación de su aplicación en el proyecto. 

\subsection{Naturaleza de la información de entrada}

La información de entrada del sistema está conformada por \textbf{imágenes de árboles recolectadas directamente en campo} en \textit{el bosque de Chapultepec en la Ciudad de México}, seleccionado como caso de estudio.

Estas imágenes constituyen el contenido del dataset, el cual representa la principal fuente de datos para \textbf{analizar visualmente }aspectos relacionados con la presencia de afecciones en el estado de salud de los árboles, provocadas por insectos descortezadores.

En este sentido, el sistema debe ser capaz de \textbf{identificar ciertos patrones} presentes en la apariencia externa de los árboles a partir de las imágenes, extrayendo características relevantes como variaciones en la textura de la corteza, alteraciones en la coloración, irregularidades superficiales o signos visibles de daño que puedan asociarse con la presencia de la plaga.

Cabe destacar que la elección de imágenes como recurso principal responde a la naturaleza misma del problema, ya que los indicios de afectación suelen manifestarse visualmente en el árbol antes de poder cuantificarse mediante otro tipo de mediciones. Por ello, el análisis computacional debe centrarse en interpretar \textbf{información visual compleja}, no estructurada y susceptible a variaciones propias del entorno de captura, tales como la iluminación, la distancia, el ángulo, el fondo y la calidad de la imagen.
